\subsection{Convergence Theorems}

% **Role:** modern integration and almost-everywhere analysis.
%
% Monotone convergence, Fatou, dominated convergence, uniform integrability preview.

This subsection collects the main limit-interchange tools: monotone convergence, Fatou, and dominated convergence.

\begin{theorem} \label{thm:monotone-convergence-integral} \hlt{Monotone Convergence}\\
Let $E\in\mathcal{M}$, let $f_n,f$ be measurable on $E$, and assume $0\leq f_n\uparrow f$. Then $\int_E f_n\,d\mu\uparrow\int_E f\,d\mu$.
\begin{proof}
By monotonicity, $\int_Ef_n\,d\mu\uparrow L\leq\int_Ef\,d\mu$.\\
For $s\in\mathcal{S}^+(f)$ and $0<c<1$, set $E_n\defeq\{x\in E:f_n(x)\geq cs(x)\}$.\\
Since $f_n\uparrow f$ and $s\leq f$, $E_n\uparrow E$, and $\int_Ef_n\,d\mu\geq\int_{E_n}f_n\,d\mu\geq c\int_{E_n}s\,d\mu$.\\
Continuity from below for $\phi_s(A)\defeq\int_As\,d\mu$ gives $L\geq c\int_Es\,d\mu$. Let $c\uparrow1$, then take $\sup\limits_{s\in\mathcal{S}^+(f)}$.\\
\hlt{Structural Commentary}: Nonnegative integration commutes with monotone limits.\\
\hlt{Key Insight}: Simple minorants of the limit are eventually captured below the sequence.
\end{proof}
\end{theorem}

\begin{theorem} \label{thm:nonnegative-series-integral} \hlt{Series Form of Monotone Convergence}\\
Let $E\in\mathcal{M}$, let $f_n\geq0$ be measurable, and set $f\defeq\sum\limits_{n=1}^{\infty}f_n$ on $E$. Then $\int_E f\,d\mu=\sum\limits_{n=1}^{\infty}\int_E f_n\,d\mu$.
\begin{proof} Set $S_N\defeq\sum\limits_{n=1}^{N}f_n$. Then $0\leq S_N\uparrow f$.\\
By Theorem \ref{thm:monotone-convergence-integral}, $\int_Ef\,d\mu=\lim\limits_{N\to\infty}\int_ES_N\,d\mu =\lim\limits_{N\to\infty}\sum\limits_{n=1}^{N}\int_Ef_n\,d\mu =\sum\limits_{n=1}^{\infty}\int_Ef_n\,d\mu$.\\
\hlt{Structural Commentary}: Nonnegative series are monotone partial sums.\\
\hlt{Key Insight}: Integrate partial sums, then pass upward.
\end{proof}
\end{theorem}

\begin{corollary} \label{cor:beppo-levi-bounded-integral-criterion} \hlt{Beppo Levi Bounded-Integral Criterion}\\
Let $E\in\mathcal{M}$, let $(f_n)$ be measurable on $E$, and assume $0\leq f_n\uparrow f$ with $\sup\limits_n\int_Ef_n\,d\mu<+\infty$.\\
Then $f<+\infty$ a.e. on $E$ and $\int_Ef\,d\mu=\lim\limits_n\int_Ef_n\,d\mu<+\infty$.
\begin{proof}
The monotone limit $f$ is measurable, and
$\int_Ef\,d\mu=\lim\limits_n\int_Ef_n\,d\mu\leq\sup\limits_n\int_Ef_n\,d\mu<+\infty$ by Theorem \ref{thm:monotone-convergence-integral}. Thus $f<+\infty$ a.e. on $E$ by Corollary \ref{cor:integral-null-and-finite-ae-consequences}.\\
\hlt{Structural Commentary}: Bounded mass survives monotone limits.\\
\hlt{Key Insight}: Finite nonnegative integral forces a.e. finiteness.
\end{proof}
\end{corollary}

\begin{theorem} \label{thm:fatou-lemma} \hlt{Fatou's Lemma}\\
Let $E\in\mathcal{M}$ and let $f_n:E\to[0,+\infty]$ be measurable. Then $\int_E\liminf\limits_nf_n\,d\mu\leq\liminf\limits_n\int_Ef_n\,d\mu$.
\begin{proof} Set $g_n\defeq\inf\limits_{i\geq n}f_i$. Then $0\leq g_n\uparrow\liminf\limits_if_i$ and $\forall i\geq n$, $g_n\leq f_i$, so $\int_Eg_n\,d\mu\leq\inf\limits_{i\geq n}\int_Ef_i\,d\mu$.\\
By Theorem \ref{thm:monotone-convergence-integral}, $\int_E\liminf\limits_if_i\,d\mu=\lim\limits_n\int_Eg_n\,d\mu\leq\lim\limits_n\inf\limits_{i\geq n}\int_Ef_i\,d\mu=\liminf\limits_n\int_Ef_n\,d\mu$.\\
\hlt{Structural Commentary}: Replace $\liminf$ by increasing tail infima.\\
\hlt{Key Insight}: Tail infima are monotone minorants.
\end{proof}
\end{theorem}

\begin{example} \label{ex:fatou-strict-concentration-and-escape} \hlt{Strictness in Fatou's Lemma}\\
Fatou's inequality can be strict by concentration or escape of mass.
\begin{enumerate}[label=(\roman*)]
\setlength{\itemsep}{0pt}
\item \hlt{Concentration}: On $E=(0,1)$ with Lebesgue measure $m$, set $f_n\defeq n\mathbf{1}_{(0,\frac{1}{n})}$.\\
Then $f_n\to0$ pointwise and $\forall n,\ \int_Ef_n\,dm=1$.
\item \hlt{Escape}: On $E=\R$ with Lebesgue measure $m$, set $g_n\defeq\mathbf{1}_{(n,n+1)}$.\\
Then $g_n\to0$ pointwise and $\forall n,\ \int_Eg_n\,dm=1$.
\end{enumerate}
In either case, with $h_n=f_n$ or $h_n=g_n$, $\int_E\liminf\limits_nh_n\,dm=0<1=\liminf\limits_n\int_Eh_n\,dm$.
\end{example}

\begin{figure}[H]
  \centering
  \begin{tikzpicture}[
  scale=0.86,
  >=Stealth,
  every node/.style={font=\scriptsize},
  axis/.style={->, black!45, line width=0.45pt},
  spike/.style={draw=blue!70!black, fill=blue!14, line width=0.7pt},
  bump/.style={draw=green!45!black, fill=green!16, line width=0.7pt},
  guide/.style={black!16, dashed, line width=0.35pt}
]
  \begin{scope}
    \draw[axis] (-0.15,0) -- (1.9,0) node[right] {$x$};
    \draw[axis] (0,-0.12) -- (0,2.05) node[above] {$y$};

    \draw[guide] (0.68,0) -- (0.68,1.72);
    \draw[guide] (0,1.72) -- (0.68,1.72);
    \draw[spike] (0,0) rectangle (0.68,1.72);

    \node[below] at (0,0) {$0$};
    \node[below] at (0.68,0) {$1/n$};
    \node[left] at (0,1.72) {$n$};
  \end{scope}

  \begin{scope}[xshift=3.7cm]
    \draw[axis] (-0.15,0) -- (4.65,0) node[right] {$x$};
    \draw[axis] (0,-0.12) -- (0,2.05) node[above] {$y$};

    \foreach \x/\lab in {0.45/1,1.55/2,2.65/3} {
      \draw[bump] (\x,0) rectangle ++(0.62,0.92);
      \node[below] at ({\x+0.34},0) {$\lab$};
    }
    \draw[bump] (3.68,0) rectangle ++(0.62,0.92);
    \node[below] at (3.99,0) {$n$};
  \end{scope}
\end{tikzpicture}

  \caption{Strict Fatou behavior: \textcolor{blue!65!black}{$f_n=n\mathbf{1}_{(0,\frac{1}{n})}$} and \textcolor{green!45!black}{$g_n=\mathbf{1}_{(n,n+1)}$} have integral $1$ but vanish pointwise.}
  \label{fig:fatou-concentration-escape}
\end{figure}

\begin{theorem} \label{thm:dominated-convergence-lebesgue} \hlt{Dominated Convergence}\\
Let $E\in\mathcal{M}$, let $f_n,f$ be measurable on $E$, and let $g\in\mathscr{L}^1(\mu)$ on $E$.\\
Suppose $f_n\to f$ a.e. on $E$ and $\forall n,\ \abs{f_n}\leq g$ on $E$. Then:
\begin{enumerate}[label=(\roman*)]
\setlength{\itemsep}{0pt}
\item $f\in\mathscr{L}^1(\mu)$ on $E$.
\item $\lim\limits_{n\to\infty}\int_Ef_n\,d\mu=\int_Ef\,d\mu$.
\end{enumerate}
\begin{proof} Choose null $N\subseteq E$ such that $f_n\to f$ on $E\setminus N$.\\
Changing values on $N$ does not change integrals, so assume $f_n\to f$ everywhere on $E$.
\begin{enumerate}[label=(\roman*)]
\setlength{\itemsep}{0pt}
\item The pointwise limit $f$ is measurable on $E$, and $\abs f\leq g$ on $E$; hence $f\in\mathscr{L}^1(\mu)$ on $E$ by Theorem \ref{thm:absolute-integrability-and-domination}(ii).
\item Put $a_n\defeq\int_Ef_n\,d\mu$. Fatou applied to $g+f_n\geq0$ gives $\int_E(g+f)\,d\mu\leq\int_Eg\,d\mu+\liminf\limits_na_n$.\\
Fatou applied to $g-f_n\geq0$ gives $\int_E(g-f)\,d\mu\leq\int_Eg\,d\mu-\limsup\limits_na_n$.\\
Subtracting $\int_Eg\,d\mu<+\infty$ gives $\int_Ef\,d\mu\leq\liminf\limits_na_n\leq\limsup\limits_na_n\leq\int_Ef\,d\mu$.
\end{enumerate}
\hlt{Structural Commentary}: Domination reduces signed convergence to two Fatou bounds.\\
\hlt{Key Insight}: Apply Fatou to $g+f_n$ and $g-f_n$.
\end{proof}
\end{theorem}

\begin{theorem} \label{thm:general-dominated-convergence-lebesgue} \hlt{General Dominated Convergence}\\
Let $E\in\mathcal{M}$, let $f_n,f,g$ be measurable on $E$, and let $g_n\geq0$ be integrable on $E$.\\
Suppose $f_n\to f$ and $g_n\to g$ a.e. on $E$, with $\forall n,\ \abs{f_n}\leq g_n$ a.e. on $E$.\\
Assume $\lim\limits_n\int_Eg_n\,d\mu=\int_Eg\,d\mu<+\infty$. Then:
\begin{enumerate}[label=(\roman*)]
\setlength{\itemsep}{0pt}
\item $\forall n,\ f_n\in\mathscr{L}^1(\mu)$ on $E$, and $f\in\mathscr{L}^1(\mu)$ on $E$.
\item $\lim\limits_{n\to\infty}\int_Ef_n\,d\mu=\int_Ef\,d\mu$.
\end{enumerate}
\begin{proof} Choose null $N\subseteq E$ so that $f_n\to f$, $g_n\to g$, and $\forall n,\ \abs{f_n}\leq g_n$ on $E\setminus N$.\\
Changing values on $N$ does not change integrals, so assume these relations hold on $E$.
\begin{enumerate}[label=(\roman*)]
\setlength{\itemsep}{0pt}
\item Since $\abs{f_n}\leq g_n\in\mathscr{L}^1(\mu)$, $f_n\in\mathscr{L}^1(\mu)$ on $E$ by Theorem \ref{thm:absolute-integrability-and-domination}(ii).\\
Passing to the limit gives $\abs f\leq g$, and $\int_Eg\,d\mu<+\infty$, hence $f\in\mathscr{L}^1(\mu)$ on $E$.
\item Put $a_n\defeq\int_Ef_n\,d\mu$, $b_n\defeq\int_Eg_n\,d\mu$, and $b\defeq\int_Eg\,d\mu$; then $b_n\to b$.\\
Fatou applied to $g_n+f_n\geq0$ gives $b+\int_Ef\,d\mu\leq\liminf\limits_n(b_n+a_n)=b+\liminf\limits_na_n$.\\
Fatou applied to $g_n-f_n\geq0$ gives $b-\int_Ef\,d\mu\leq\liminf\limits_n(b_n-a_n)=b-\limsup\limits_na_n$.\\
Subtracting $b<+\infty$ gives $\limsup\limits_na_n\leq\int_Ef\,d\mu\leq\liminf\limits_na_n$, hence $a_n\to\int_Ef\,d\mu$.
\end{enumerate}
\hlt{Structural Commentary}: Varying domination works when masses converge.\\
\hlt{Key Insight}: Fatou on $g_n\pm f_n$ traps the limit.
\end{proof}
\end{theorem}

\begin{corollary} \label{cor:monotone-convergence-ae} \hlt{Monotone Convergence Almost Everywhere}\\
Let $E\in\mathcal{M}$ and $0\leq f_n\uparrow f$ a.e. on $E$. Then $\int_Ef_n\,d\mu\uparrow\int_Ef\,d\mu$.
\begin{proof}
Choose $N\subseteq E$ with $\mu(N)=0$ and $f_n\uparrow f$ on $E\setminus N$. Set $g_n\defeq f_n\mathbf{1}_{E\setminus N}$ and $g\defeq f\mathbf{1}_{E\setminus N}$.\\
Then $0\leq g_n\uparrow g$, so $\int_Eg_n\,d\mu\uparrow\int_Eg\,d\mu$. Since $f_n=g_n$ and $f=g$ off $N$, the integrals over $N$ vanish.\\
\hlt{Structural Commentary}: Monotone convergence ignores null exceptional sets.\\
\hlt{Key Insight}: Remove the exceptional set, prove the theorem there, then add back zero.
\end{proof}
\end{corollary}

\begin{example} \label{ex:counting-measure-integral-is-series} \hlt{Counting Measure Integrals Are Series}\\
Let $\mu$ be counting measure on $(\N,2^\N)$. For $f:\N\to[0,+\infty]$, set $s_N\defeq\sum\limits_{n=1}^{N}f(n)\mathbf{1}_{\{n\}}$; then $0\leq s_N\uparrow f$, with the case $f(k)=+\infty$ giving both sides $+\infty$. Hence $\int_{\N} f\,d\mu=\lim\limits_N\int_{\N}s_N\,d\mu =\lim\limits_N\sum\limits_{n=1}^{N}f(n) =\sum\limits_{n=1}^{\infty}f(n)$.\\
\hlt{Structural Commentary}: Counting measure turns integration into summation.\\
\hlt{Key Insight}: Finite-support truncations are the simple-function partial sums.
\end{example}

\begin{example} \label{ex:monotone-fatou-boundary-tests} \hlt{Boundary Tests for Monotone Convergence and Fatou}\\
On $([0,+\infty),\mathcal{B},m)$:
\begin{enumerate}[label=(\roman*)]
\setlength{\itemsep}{0pt}
\item $f_n\defeq\mathbf{1}_{[0,n]}\uparrow\mathbf{1}_{[0,+\infty)}$ and $\int f_n\,dm=n\uparrow+\infty$; monotone convergence allows $+\infty$.
\item $g_n\defeq n^{-1}\mathbf{1}_{[n,+\infty)}\downarrow0$ uniformly, but $\int g_n\,dm=+\infty$ $\forall n$; decreasing convergence needs finite mass.
\item $w_n\defeq-n^{-1}\mathbf{1}_{[0,n]}\rightrightarrows0$, but $\int w_n\,dm=-1$; Fatou can fail without $f_n\geq0$ or $f_n\geq-h\in L^1$.
\end{enumerate}
\end{example}

\begin{theorem} \label{thm:fatou-with-integrable-lower-bound} \hlt{Fatou with an Integrable Lower Bound}\\
Let $h\geq0$, $h\in\mathscr{L}^1(\mu)$, and let $f_n$ be measurable with $f_n\geq-h$.\\
Then, whenever the signed integrals are defined in $\overline{\R}$, $\int \liminf\limits_{n\to\infty} f_n\,d\mu\leq\liminf\limits_{n\to\infty}\int f_n\,d\mu$.
\begin{proof}
Set $F_n\defeq f_n+h\geq0$. Fatou gives $\int \liminf\limits_nF_n\,d\mu\leq\liminf\limits_n\int F_n\,d\mu$.\\
Since $\liminf\limits_nF_n=h+\liminf\limits_nf_n$ and $\int F_n\,d\mu=\int f_n\,d\mu+\int h\,d\mu$, subtract the finite term $\int h\,d\mu$.\\
\hlt{Structural Commentary}: Fatou survives controlled negative parts.\\
\hlt{Key Insight}: Add an integrable floor to make the sequence nonnegative.
\end{proof}
\end{theorem}

\begin{remark} \label{rem:negative-fatou-example-no-integrable-floor} \hlt{Why the Negative Example Has No Floor}\\
For $w_n\defeq-n^{-1}\mathbf{1}_{[0,n]}$, any $h\in L^1$ with $-h\leq w_n$ $\forall n$ must satisfy $h\geq n^{-1}$ on $[0,n]$.\\
Thus $h\geq k^{-1}$ on $[k-1,k]$, and $\int_0^N h\,dm\geq H_N\defeq\sum\limits_{k=1}^{N}\frac{1}{k}\to+\infty$, contradicting $h\in L^1$.
\end{remark}

\begin{perspective} \label{per:domination-as-limit-integral-safeguard} \hlt{Domination Is the Safeguard}\\
$f_n\to f$ a.e. need not imply $\int f_n\to\int f$; domination $\abs{f_n}\leq g\in L^1$ restores the limit interchange.
\end{perspective}

\begin{example} \label{ex:dominated-convergence-improper-integral} \hlt{A Nonuniform Limit with Vanishing Integrals}\\
For $x>0$, set $f_n(x)\defeq e^{-nx}x^{-\frac{1}{2}}$. Then $f_n\to0$ pointwise, but $\forall\delta>0$, $\sup\limits_{0<x<\delta}f_n(x)=+\infty$.\\
For $n\geq1$, $0\leq f_n\leq g\defeq e^{-x}x^{-\frac{1}{2}}$; $g\leq x^{-\frac{1}{2}}\mathbf{1}_{(0,1]}+e^{-x}\mathbf{1}_{[1,\infty)}\in\mathscr{L}^1(0,\infty)$.\\
Thus dominated convergence gives $I_n\defeq\int_0^\infty e^{-nx}x^{-\frac{1}{2}}\,dx\to0$.\\
\hlt{Structural Commentary}: Uniform convergence is unnecessary once one integrable envelope controls the sequence.\\
\hlt{Key Insight}: Domination handles the singular behavior at $0$ and the tail at infinity simultaneously.
\end{example}

\begin{proposition} \label{prop:uniform-convergence-finite-measure-integrals} \hlt{Uniform Convergence on Finite Measure Spaces}\\
Assume $\mu(X)<+\infty$.
\begin{enumerate}[label=(\roman*)]
\setlength{\itemsep}{0pt}
\item If $0\leq f_n$ are measurable and $f_n\rightrightarrows f$, then $f\geq0$ is measurable and $\int_X f\,d\mu=\lim\limits_n\int_X f_n\,d\mu$ in $[0,+\infty]$.
\item If $f_n\in\mathscr{L}^1(\mu)$ and $f_n\rightrightarrows f$, then $f\in\mathscr{L}^1(\mu)$ and $\int_Xf_n\,d\mu\to\int_Xf\,d\mu$ in $\R$.
\end{enumerate}
\begin{proof}
\begin{enumerate}[label=(\roman*)]
\setlength{\itemsep}{0pt}
\item Measurability and $f\geq0$ pass to pointwise limits. For $\varepsilon>0$ and $n\gg1$, $f\leq f_n+\varepsilon\mathbf{1}_X$ and $f_n\leq f+\varepsilon\mathbf{1}_X$.\\
Thus $\int_X f\,d\mu\leq \int_X f_n\,d\mu+\varepsilon\mu(X)$ and $\int_X f_n\,d\mu\leq \int_X f\,d\mu+\varepsilon\mu(X)$. Let $\varepsilon\downarrow0$.
\item Choose $N$ with $\abs{f-f_N}\leq1$. Then $f=f_N+(f-f_N)\in\mathscr{L}^1(\mu)$, and\\
$\abs{\int_Xf_n\,d\mu-\int_Xf\,d\mu}\leq\int_X\abs{f_n-f}\,d\mu\leq\mu(X)\sup\limits_X\abs{f_n-f}\to0$.
\end{enumerate}
\hlt{Structural Commentary}: On finite measure spaces, uniform error controls integral error.\\
\hlt{Key Insight}: Finiteness of $\mu(X)$ turns a supremum bound into an integral bound.
\end{proof}
\end{proposition}

\begin{theorem} \label{thm:bounded-convergence-finite-measure} \hlt{Bounded Convergence on Finite Measure Spaces}\\
Let $\mu(X)<+\infty$, and let $f_n,f:X\to\R$ be measurable.\\
Suppose $f_n\to f$ a.e. and $\exists M<+\infty$ such that $\forall n,\ \abs{f_n}\leq M$ a.e. Then:
\begin{enumerate}[label=(\roman*)]
\setlength{\itemsep}{0pt}
\item $f\in\mathscr{L}^1(\mu)$.
\item $\lim\limits_{n\to\infty}\int_Xf_n\,d\mu=\int_Xf\,d\mu$.
\end{enumerate}
\begin{proof} Since $\abs{f_n}\leq M$ a.e. and $f_n\to f$ a.e., $\abs f\leq M$ a.e.\\
Set $G\defeq M\mathbf{1}_X$. Then $G\in\mathscr{L}^1(\mu)$, since $\int_XG\,d\mu=M\mu(X)<+\infty$.
\begin{enumerate}[label=(\roman*)]
\setlength{\itemsep}{0pt}
\item Since $\abs f\leq G$ a.e., $f\in\mathscr{L}^1(\mu)$ by Theorem \ref{thm:absolute-integrability-and-domination}(ii).
\item Since $\forall n,\ \abs{f_n}\leq G$ a.e., $\int_Xf_n\,d\mu\to\int_Xf\,d\mu$ by Theorem \ref{thm:dominated-convergence-lebesgue}(ii).
\end{enumerate}
\hlt{Structural Commentary}: On a finite-measure space, a uniform pointwise bound is an integrable dominator.\\
\hlt{Key Insight}: Bounded convergence is the finite-measure special case of dominated convergence.
\end{proof}
\end{theorem}

\begin{example} \label{ex:uniform-convergence-needs-finite-measure} \hlt{Why the Finite-Measure Hypothesis Matters}\\
On $(\R,m)$, let $f_n\defeq n^{-1}\mathbf{1}_{[0,n]}$. Then $\norm{f_n}_\infty=\frac{1}{n}\to0$, but $\int_\R f_n\,dm=1$ $\forall n$.\\
Thus uniform convergence need not commute with integrals when $m(\R)=+\infty$; Fatou only yields $0\leq1$.
\end{example}

\begin{example} \label{ex:dominated-convergence-needs-domination} \hlt{Why Domination Cannot Be Dropped}\\
On $(\R,m)$, let $f_n\defeq n\mathbf{1}_{[\frac{1}{n},\frac{2}{n}]}$. Then $f_n\to0$ pointwise, but $\int_\R f_n\,dm=1$ $\forall n$.\\
Thus pointwise convergence without domination need not commute with integrals; Fatou only yields $0\leq1$.
\end{example}

\begin{theorem} \label{thm:l1-absolutely-summable-series-integrate-termwise} \hlt{Absolutely Summable $L^1$ Series}\\
Let $f_n\in\mathscr{L}^1(\mu)$ and $\sum\limits_{n\geq1}\norm{f_n}_1<+\infty$.\\
Set $S_N\defeq\sum\limits_{n=1}^{N}f_n$. Then $\exists f:X\to\R$ measurable such that:
\begin{enumerate}[label=(\roman*)]
\setlength{\itemsep}{0pt}
\item $S_N\to f$ absolutely a.e.
\item $f\in\mathscr{L}^1(\mu)$.
\item $\int_X f\,d\mu=\sum\limits_{n\geq1}\int_X f_n\,d\mu$.
\end{enumerate}
\begin{proof}
Set $G\defeq\sum\limits_{n\geq1}\abs{f_n}$. By monotone convergence, $\int_XG\,d\mu=\sum\limits_{n\geq1}\norm{f_n}_1<+\infty$, so $G<+\infty$ a.e. and $G\in\mathscr{L}^1(\mu)$. Define $f\defeq\lim\limits_NS_N$ on $\{G<+\infty\}$ and $f\defeq0$ elsewhere.
\begin{enumerate}[label=(\roman*)]
\setlength{\itemsep}{0pt}
\item Since $G<+\infty$ a.e., $\sum\limits_{n\geq1}\abs{f_n}<+\infty$ a.e.; hence $S_N\to f$ absolutely a.e.
\item Since $\abs f\leq G$ a.e. and $G\in\mathscr{L}^1(\mu)$, $f\in\mathscr{L}^1(\mu)$ by Theorem \ref{thm:absolute-integrability-and-domination}(ii).
\item Since $\abs{S_N}\leq G$ and $S_N\to f$ a.e., dominated convergence gives\\
$\int_Xf\,d\mu=\lim\limits_N\int_XS_N\,d\mu=\lim\limits_N\sum\limits_{n=1}^{N}\int_Xf_n\,d\mu=\sum\limits_{n\geq1}\int_Xf_n\,d\mu$.
\end{enumerate}
\hlt{Structural Commentary}: Absolute summability in $L^1$ gives a.e. absolute convergence and termwise integration.\\
\hlt{Key Insight}: First integrate the nonnegative majorant $\sum\abs{f_n}$.
\end{proof}
\end{theorem}

\begin{proposition} \label{prop:l1-convergence-controls-absolute-integrals} \hlt{$L^1$ Convergence Controls Absolute Integrals}\\
If $f_n,f\in\mathscr{L}^1(\mu)$ and $\norm{f_n-f}_1\to0$, then:
\begin{enumerate}[label=(\roman*)]
\setlength{\itemsep}{0pt}
\item $\norm{f_n}_1\to\norm{f}_1$.
\item $\int_X\abs{f}\,d\mu=\lim\limits_n\int_X\abs{f_n}\,d\mu$.
\end{enumerate}
\begin{proof}
\begin{enumerate}[label=(\roman*)]
\setlength{\itemsep}{0pt}
\item The reverse triangle inequality gives $\abs{\abs{f_n}-\abs{f}}\leq\abs{f_n-f}$; hence $\abs{\norm{f_n}_1-\norm{f}_1}\leq\norm{f_n-f}_1\to0$.
\item Since $\norm{h}_1=\int_X\abs h\,d\mu$, item (i) gives $\int_X\abs{f_n}\,d\mu\to\int_X\abs f\,d\mu$.
\end{enumerate}
\hlt{Structural Commentary}: $L^1$ convergence is strong enough to converge the $L^1$ sizes.\\
\hlt{Key Insight}: The reverse triangle inequality converts function convergence into norm convergence.
\end{proof}
\end{proposition}

\begin{proposition} \label{prop:finite-positive-integral-localization} \hlt{Finite Integral Localizes to Finite Measure}\\
Let $f:X\to[0,+\infty]$ be measurable and $\int_X f\,d\mu<+\infty$. For $n\geq1$, set $E_n\defeq\{\frac{1}{n}<f\leq n\}$. Then:
\begin{enumerate}[label=(\roman*)]
\setlength{\itemsep}{0pt}
\item $\forall n,\ \mu(E_n)<+\infty$.
\item $\int_{X\setminus E_n}f\,d\mu\downarrow0$.
\end{enumerate}
Consequently, $\forall\varepsilon>0\ \exists E\in\mathcal{M}$ with $\mu(E)<+\infty$ and $\int_{X\setminus E}f\,d\mu<\varepsilon$.
\begin{proof}
\begin{enumerate}[label=(\roman*)]
\setlength{\itemsep}{0pt}
\item Since $f>\frac{1}{n}$ on $E_n$, $\mu(E_n)\leq n\int_{E_n}f\,d\mu\leq n\int_Xf\,d\mu<+\infty$.
\item Since $f\mathbf{1}_{E_n}\uparrow f$ a.e., monotone convergence gives $\int_{E_n}f\,d\mu\uparrow\int_Xf\,d\mu$.\\
Thus $\int_{X\setminus E_n}f\,d\mu=\int_Xf\,d\mu-\int_{E_n}f\,d\mu\downarrow0$.
\end{enumerate}
Given $\varepsilon>0$, choose $n$ with $\int_{X\setminus E_n}f\,d\mu<\varepsilon$ and set $E=E_n$.\\
\hlt{Structural Commentary}: Finite integral mass can be captured on a finite-measure part of the space.\\
\hlt{Key Insight}: Keep the levels where $f$ is bounded above and away from zero.
\end{proof}
\end{proposition}

\begin{proposition} \label{prop:setwise-convergence-from-total-mass-convergence} \hlt{Setwise Convergence from Total Mass Convergence}\\
Let $f_n,f:X\to[0,+\infty]$ be measurable, $f_n\to f$ pointwise, and $\int_X f\,d\mu=\lim\limits_{n\to\infty}\int_X f_n\,d\mu<+\infty$.\\
Then $\forall E\in\mathcal{M}$, $\int_E f\,d\mu=\lim\limits_{n\to\infty}\int_E f_n\,d\mu$.
\begin{proof}
Fix $E\in\mathcal{M}$ and set $E^c\defeq X\setminus E$. By Fatou, $\int_Ef\,d\mu\leq\liminf\limits_n\int_Ef_n\,d\mu$ and $\int_{E^c}f\,d\mu\leq\liminf\limits_n\int_{E^c}f_n\,d\mu$.\\
Since $\int_X f_n\,d\mu=\int_E f_n\,d\mu+\int_{E^c}f_n\,d\mu\to\int_Xf\,d\mu<+\infty$, $\limsup\limits_n\int_Ef_n\,d\mu\leq\int_Xf\,d\mu-\int_{E^c}f\,d\mu=\int_Ef\,d\mu$.\\
Thus $\liminf\limits_n\int_Ef_n\,d\mu=\limsup\limits_n\int_Ef_n\,d\mu=\int_Ef\,d\mu$.\\
\hlt{Structural Commentary}: Total mass plus Fatou on complements fixes each part.\\
\hlt{Key Insight}: A lower bound on the complement is an upper bound on the chosen set.
\end{proof}
\end{proposition}

\begin{definition} \label{def:function-truncates} \hlt{Function Truncates}\\
For $f:X\to\R$ and $n\in\N$, define the \hlt{$n$th truncate}
\begin{equation*}
T_nf(x)\defeq (-n)\vee(f(x)\wedge n)=
\begin{cases}
-n,&f(x)<-n,\\
f(x),&-n\leq f(x)\leq n,\\
n,&f(x)>n.
\end{cases}
\end{equation*}
\end{definition}

\begin{proposition} \label{prop:truncation-characterization-integrability} \hlt{Truncates Characterize Integrability}\\
Let $f:X\to\R$ be measurable and set $f_n\defeq T_nf$ for $n\in\N$.
\begin{enumerate}[label=(\roman*)]
\setlength{\itemsep}{0pt}
\item $f\in\mathscr{L}^1(\mu)\Longleftrightarrow \sup_n\norm{f_n}_1<+\infty$.
\item If $f\in\mathscr{L}^1(\mu)$, then $\int_Xf_n\,d\mu\to\int_Xf\,d\mu$.
\end{enumerate}
\begin{proof} $\forall n$, $T_n:\R\to[-n,n]$ is continuous, so $f_n$ is measurable; moreover $f_n\to f$, $\abs{f_n}\leq\abs f$, and $\abs{f_n}\uparrow\abs f$.
\begin{enumerate}[label=(\roman*)]
\setlength{\itemsep}{0pt}
\item[$(i)$]$(\Rightarrow)$ If $f\in\mathscr{L}^1(\mu)$, then $\norm{f_n}_1\leq\norm{f}_1$ for all $n$, so $\sup_n\norm{f_n}_1<+\infty$. 
\item[$(i)$]$(\Leftarrow)$ If $\sup_n\norm{f_n}_1<+\infty$, then MCT gives $\norm{f}_1=\lim_{n\to\infty}\norm{f_n}_1=\sup_n\norm{f_n}_1<+\infty$, so $f\in\mathscr{L}^1(\mu)$.
\item[$(ii)$] If $f\in\mathscr{L}^1(\mu)$, then DCT applied to $f_n\to f$ with $\abs{f_n}\leq\abs f$ gives $\int_Xf_n\,d\mu\to\int_Xf\,d\mu$.
\end{enumerate}
\hlt{Structural Commentary}: Truncation reduces $L^1$ questions to bounded approximants.\\
\hlt{Key Insight}: MCT reads $\norm{f}_1$ from $\abs{f_n}\uparrow\abs f$; DCT passes signed integrals.
\end{proof}
\end{proposition}

\begin{example} \label{ex:setwise-convergence-needs-finite-total-mass} \hlt{Why the Finite Total Mass Assumption Matters}\\
Let $X=\N\times\{0,1\}$ with counting measure, set $E\defeq\N\times\{1\}$.\\
Define $f\defeq\mathbf{1}_{\N\times\{0\}}$; $f_n\defeq\mathbf{1}_{\{1,\ldots,n\}\times\{0\}}+n\mathbf{1}_{\{n\}\times\{1\}}$.\\
Then $f_n\to f$ pointwise and $\int_Xf\,d\mu=\lim\limits_n\int_Xf_n\,d\mu=+\infty$; $\int_Ef\,d\mu=0,\; \int_Ef_n\,d\mu=n$.\\
Thus Proposition \ref{prop:setwise-convergence-from-total-mass-convergence} can fail when the common total mass is infinite.
\end{example}
